\documentclass{article}

% If you need to pass options to natbib, use, e.g.:
%     \PassOptionsToPackage{numbers, compress}{natbib}
% before loading neurips_2026.

\usepackage[eandd]{neurips_2026}

\usepackage[utf8]{inputenc}
\usepackage[T1]{fontenc}
\usepackage{hyperref}
\usepackage{url}
\usepackage{booktabs}
\usepackage{amsfonts}
\usepackage{amsmath}
\usepackage{amssymb}
\usepackage{nicefrac}
\usepackage{microtype}
\usepackage{xcolor}
\usepackage{graphicx}
\usepackage{array}

\newcommand{\meg}{\mathrm{MEG}}
\newcommand{\mig}{\mathrm{MIG}}
\newcommand{\qbench}{Q_{\text{bench}}}

\title{\textsc{Veil}: A Hidden-Evaluation Benchmark for Multi-Agent Structured Optimization}

\author{%
  Anonymous Author(s)
}

\begin{document}

\maketitle

\begin{abstract}
Different models bring different strengths, so combining them should help. But to combine them, multi-agent protocols make them communicate---and communication makes them more alike. We call this the \emph{diversity tax}: the cost that interaction imposes on the independence it needs.

On nine structured optimization tasks scored by hidden verifiers that agents never observe, we find that model-family diversity is the only factor that helps ($+0.19$, 95\% CI $[+0.06, +0.31]$; $N{=}120$ balanced factorial). Every interaction mechanism we test---debate, iterative refinement, hierarchical planning---pays the diversity tax in this benchmark: it pushes diverse agents toward consensus and erases their complementary information. Only one-shot composition, where no agent sees another's work, preserves independence. The distinction matters: \emph{composition is not interaction}. Across four model families and 1{,}852 runs, no multi-agent protocol robustly exceeds matched single-agent controls. Doubling the budget does not change this.
\end{abstract}

% ══════════════════════════════════════════════════════════════════════════════
\section{Introduction}
\label{sec:intro}

The premise of multi-agent LLM systems is that combining different models should be better than using one. Different model families encode different training distributions, make different errors, and find different solutions. In ensemble learning, this independence is the entire source of value: the ambiguity decomposition \citep{krogh1994neural} shows that ensemble accuracy depends on member diversity, and adding a correlated member can make things worse. But every multi-agent protocol in current use---debate, iterative refinement, hierarchical planning---works by making agents communicate. Communication creates correlation. This is the puzzle: \emph{the mechanism designed to exploit diversity is the mechanism that destroys it}.

We call this the \emph{diversity tax}. On nine structured optimization tasks scored by hidden verifiers that agents never observe, we find that model-family diversity is the only factor that significantly improves multi-agent performance ($+0.19$, CI $[+0.06, +0.31]$; balanced $2{\times}2$ factorial, $N{=}120$). But every protocol that makes agents communicate pays the tax: debate, iterative refinement, and hierarchical planning all push diverse agents toward consensus and erase their complementary information. Only one-shot composition---where a synthesizer reads three independent proposals and no proposer ever sees another's work---preserves independence. \emph{Composition is not interaction.} This distinction, invisible when all agents share a model family, is the central empirical finding of this benchmark. Across four model families and 1{,}852 runs, no multi-agent protocol robustly exceeds matched single-agent controls.

Two methodological choices make this identification possible. \textbf{Hidden evaluation}: each task has a visible development evaluator for search and a held-out hidden evaluator for final scoring, preventing protocols from optimizing the scorer that judges them. \textbf{Matched-budget controls}: every protocol runs under the same resource cap ($T{=}200$K tokens, $W{=}600$s, $K{=}25$ evaluator calls) and is compared against three single-agent baselines---Self-Refine (iterative depth), Best-of-$N$ (parallel breadth), and verifier-guided search (feedback exploitation)---so gains cannot come from spending more. Budget matching is on caps, not realized usage; per-protocol consumption is reported in Appendix~\ref{app:budget}.

The contributions are:

\begin{enumerate}
\item \textbf{\textsc{Veil} benchmark}: nine structured optimization tasks with deterministic hidden evaluators, matched-budget controls, and two metrics---MEG (does a protocol beat the best single-agent control?) and MIG (does the interaction step itself add or destroy value?)---that together provide a reusable testbed for multi-agent evaluation under hidden scoring.
\item \textbf{First benchmark-scale evidence for the diversity tax}: across 1{,}852 runs and four model families, a balanced $2{\times}2$ factorial on \textsc{Veil} isolates model-family diversity as the only significant factor ($+0.19$, CI $[+0.06, +0.31]$); MIG decomposition shows that every tested interaction protocol erodes it. No multi-agent protocol robustly exceeds matched single-agent controls.
\item \textbf{A composition--interaction distinction}: on \textsc{Veil}, one-shot synthesis is the only mixed-model protocol with positive MIG, because no proposer sees another's work. The benchmark separates protocols that combine independent outputs from those that make agents communicate, and the two behave differently.
\end{enumerate}

\paragraph{Scope of claims.}
The diversity finding rests on two co-primary analyses: the $2{\times}2$ ablation ($N{=}120$, pre-specified factorial, bootstrap inference) identifies model-family diversity as the signal; the MIG decomposition identifies interaction as the mechanism that erodes it. The effect is clear on two of three ablation tasks (DiffBases, Erd\H{o}s), absent on the third (MolQED), and untestable on six low-spread tasks. The broader sweep across 41 variants (Appendix~\ref{app:allvariants}) is exploratory. All protocols are our reimplementations; we match published designs but have not validated against original code.

% ══════════════════════════════════════════════════════════════════════════════
\section{Related Work}
\label{sec:related}

\paragraph{Multi-agent methods.}
Self-Refine provides a canonical iterative single-agent baseline \citep{madaan2023selfrefine}. Multi-agent debate \citep{du2024debate}, MAgICoRe \citep{chen2025magicore}, and Mixture-of-Agents \citep{wang2025moa} supply interaction baselines spanning adversarial critique, role-structured refinement, and parallel synthesis. CAMEL \citep{li2023camel}, AutoGen \citep{wu2023autogen}, AgentVerse \citep{chen2023agentverse}, and MetaGPT \citep{hong2023metagpt} provide collaboration architectures; Plan-and-Solve \citep{wang2023plansolve} formalizes hierarchical decomposition. All prior frameworks are evaluated on tasks where the test-time evaluator is either public or an LLM judge; none use a hidden deterministic verifier that agents never observe.

\paragraph{Matched-compute controls and test-time scaling.}
Prior work under matched compute shows that strong single-agent controls erase many apparent multi-agent gains on reasoning and math tasks \citep{xu2026single,tran2026equal}. Those studies use public evaluators (exact-match correctness); our work extends the finding to optimization tasks where agents iteratively search using a visible scorer, introducing the hidden-evaluation split to detect when gains are scorer-specific artifacts. The test-time compute scaling literature \citep{snell2024scaling} frames Best-of-$N$ selection as a canonical scaling strategy. FunSearch \citep{romeraparedes2024funsearch}, AlphaEvolve \citep{novikov2025alphaevolve}, and its mathematical applications \citep{georgiev2025alphaevolve} demonstrate evolutionary search over solution artifacts. Our contribution relative to all of the above is the hidden-evaluation methodology and the identification of the diversity tax as the mechanism behind the negative finding.

\paragraph{Ensemble diversity and social influence.}
The ambiguity decomposition \citep{krogh1994neural} shows that ensemble error equals average individual error minus a diversity term; adding correlated members reduces this diversity bonus. The diversity tax we identify is a structural instance of this principle for LLM agents, where interaction creates the correlation. \citet{lorenz2011social} demonstrate the same phenomenon in human groups: social influence reduces estimate diversity and degrades collective accuracy. Multi-agent debate is, in this framing, social influence for LLMs.

% ══════════════════════════════════════════════════════════════════════════════
\section{Testbed Design}
\label{sec:testbed}

\subsection{Marginal Epistemic Gain}
\label{sec:meg}

Let $\qbench(y)$ denote the task-normalized hidden-evaluation score of a final artifact $y$, and let $B=(T,W,C,K)$ denote the observable resource budget: billed-token cap $T$, wall-clock limit $W$, development-evaluator CPU budget $C$, and evaluator call cap $K$. A run is feasible only if each component stays within its cap; the components are not exchangeable. For a protocol $p$,
\begin{equation}
\meg(p,B)=\mathbb{E}\!\left[\qbench(y_p)\right]-
\max\!\left(
\mathbb{E}\!\left[\qbench(y_{\mathrm{SR}})\right],
\mathbb{E}\!\left[\qbench(y_{\mathrm{BoN}})\right],
\mathbb{E}\!\left[\qbench(y_{\mathrm{VGS}})\right]
\right).
\label{eq:meg}
\end{equation}
The three controls span the axes a single agent can exploit: Self-Refine captures \emph{iterative depth}, Best-of-$N$ captures \emph{parallel breadth}, and VGS captures \emph{feedback-driven exploration} through a FunSearch-style evolutionary loop \citep{romeraparedes2024funsearch} with OPRO-style scored-example mutation \citep{yang2024opro}. Positive MEG therefore means value that is not reducible to deeper refinement, wider sampling, or more aggressive use of visible feedback. The binding control varies by task.

MEG is computed per task and aggregated as the unweighted mean across tasks; bootstrap CIs resample seeds within tasks and then average. Because the denominator takes the max of three sample means (each from $n{=}5$ seeds), it has a small upward bias (mean ${\approx}{+}0.006$ across tasks, though one task shows a negligible $-0.002$), pushing MEG conservatively negative on aggregate (Appendix~\ref{app:megbias}). A bootstrap bias correction shifts all protocols by $+0.006$ AggMEG with no sign changes; the ``consistent with null'' interpretation for near-zero protocols accounts for this via the CI spanning zero.

MEG answers the \emph{practitioner} question: should I switch from the best single-agent control? A complementary metric isolates whether the interaction step itself adds value. The \textbf{Marginal Interaction Gain (MIG)} compares the same agent team with and without interaction:
\begin{equation}
\mig(p,B)=\mathbb{E}\!\left[\qbench(y_p^{\,\text{interact}})\right]-
\mathbb{E}\!\left[\qbench(y_p^{\,\text{no-interact}})\right],
\label{eq:mig}
\end{equation}
where $y_p^{\,\text{interact}}$ and $y_p^{\,\text{no-interact}}$ are artifacts from the same set of agents \emph{with} and \emph{without} the interaction step, holding model-family assignment constant. For MoA, the ablation is exact: synthesis versus best-of-proposals from the same three models. For other protocol families (debate, MAgICoRe, HPE), we define $y_p^{\,\text{no-interact}}$ as the best single-agent output from the same model-family set, namely what the strongest individual model would produce without any interaction. MIG can be positive (interaction creates value) even when MEG is negative (the value does not exceed what the best single-agent control achieves independently).

\paragraph{Why this decomposition.}
MEG uses the max of three controls because a practitioner choosing a protocol would compare against the \emph{best} available single-agent alternative, not the average one. The three controls are chosen to cover the degrees of freedom a single agent can exploit---depth (Self-Refine), breadth (Best-of-$N$), and feedback exploitation (VGS)---so that positive MEG cannot be explained by any of these axes alone. The $2{\times}2$ factorial isolates diversity and synthesis because these are the two design choices a practitioner actually controls when configuring a multi-agent system: which models to include, and how to combine their outputs. MIG complements MEG by holding the team constant and toggling only the interaction step, separating the value of \emph{having} diverse agents from the value of \emph{making them communicate}. Hidden evaluation follows the same logic as a held-out test set in supervised learning: if agents can observe and optimize the scorer, gains may reflect scorer exploitation rather than solution quality. The dev/hidden split makes this testable.

\subsection{Task Suite}
\label{sec:task-suite}

The testbed targets structured artifacts scored by deterministic oracles (verifier robustness audit in Appendix~\ref{app:verifier-audit}). Each task satisfies four criteria: broad strategy space, deterministic held-out evaluation, meaningful iterative improvement under bounded feedback, and resource accounting that can be matched across protocols.

\begin{table}[t]
\caption{Task suite summary. The labeled tasks support the $2{\times}2$ ablation; TSP-50 is a discriminative task with calibrated anchors (scaled from TSP-100); MaxCut is a hidden-evaluation diagnostic; Circle Packing is a control case. The remaining tasks have low score spread under the current budget.}
\label{tab:coretracks}
\centering
\small
\begin{tabular}{l l l l}
\toprule
Task & Domain & Hidden type & Role \\
\midrule
Difference Bases & Number theory & I (95\% coverage) & $2{\times}2$ ablation \\
Erd\H{o}s overlap & Combinatorics & II (fine mesh) & $2{\times}2$ ablation \\
Molecule QED & Chemistry & II (QED$\times$SA) & $2{\times}2$ ablation \\
TSP-50 & Combinatorics & II (perturbed cities) & Discriminative \\
MaxCut $G(200,0.3)$ & Graph theory & I (full graph) & Hidden-eval diagnostic \\
Circle Packing $n{=}20$ & Geometry & I (strict tol.) & Neg.\ control$^\ddagger$ \\
Flat Poly deg.~50 & Analysis & II (adversarial pts.) & Low-variation \\
TSP-100 & Combinatorics & II (perturbed cities) & Low-variation \\
LJ $n{=}41$ clusters & Chemistry & II (3-body corr.) & Low-variation \\
\bottomrule
\end{tabular}

\smallskip\noindent\small $^\ddagger$Control condition for diversity: Circle Packing's global overlap constraint (low composability) means diversity should \emph{not} help; it confirms that the diversity effect is task-dependent, not universal.
\end{table}

Each task uses one of two hidden-evaluator designs: \textbf{Type~I} replaces the development evaluator's input (e.g., the 80\%-edge subgraph becomes the full graph for MaxCut; the loose circle tolerance becomes strict), while \textbf{Type~II} adds a new scoring component invisible during search (e.g., perturbed city coordinates for TSP, an adversarial point set for Flat Polynomials, a synthesizability penalty for MolQED). The parenthetical in Table~\ref{tab:coretracks} summarizes the specific hidden modification.

\textsc{Veil} intentionally contains four task roles, each serving a distinct benchmarking function (Table~\ref{tab:coretracks}): \emph{discriminative tasks} (DiffBases, Erd\H{o}s, MolQED, TSP-50) produce enough score variation to separate protocols and support the $2{\times}2$ ablation; a \emph{diagnostic task} (MaxCut) tests hidden-evaluation methodology; a \emph{negative control} (Circle Packing) confirms the diversity effect is task-dependent; and \emph{capability-frontier tasks} (Flat Poly, TSP-100, LJ) measure whether any protocol makes meaningful progress under the $T{=}200$K budget. Low-spread tasks are not failed tasks---they define the current capability boundary and may become discriminative as models or budgets improve. Including them is conservative: they dilute the aggregate toward null, making positive MEG harder to achieve, and they prevent the benchmark from appearing cherry-picked for discrimination. The suite still leans toward combinatorial mathematics; broader scientific coverage remains a limitation.

\paragraph{Task inclusion protocol.}
A task enters the testbed only if it has a stable verifier, a non-saturated baseline, non-trivial search headroom, and frozen normalization anchors (Appendix~\ref{app:normalization}).

\paragraph{Benchmark interface.}
\textsc{Veil} is released as a self-contained artifact: frozen task instances (problem descriptions, solution schemas, and parameters), development evaluators (visible during search), hidden evaluators (used only for final scoring and never exposed to agents), prompt templates for all protocols, a protocol harness that enforces budget caps and logs token/call traces, and a result cache of all runs reported in this paper. To evaluate a new protocol, a user implements a single function mapping (task, budget) $\to$ artifact; the harness handles evaluation, budget enforcement, submission, and MEG computation. New tasks can be added by supplying a dev evaluator, a hidden evaluator, and calibrated normalization anchors following the procedure in Appendix~\ref{app:normalization}. The extensibility requirements and hidden-evaluator design patterns are detailed in Appendix~\ref{app:extensibility}.

\paragraph{Testbed versioning.}
\textsc{Veil} v1.0 pins all model identifiers, prompt templates, task instances, and evaluators at freeze time. One mid-run model substitution occurred: the Gemini preview model was decommissioned by OpenRouter and replaced with the GA release; all affected cells were re-run.

\subsection{Protocol Set}
\label{sec:protocols}

The primary comparison set contains ten protocols under a matched per-task budget ($T{=}200$K tokens, $W{=}600$s wall-clock, $C{=}30$s evaluator CPU, $K{=}25$ evaluator calls): single-shot, Self-Refine \citep{madaan2023selfrefine}, Best-of-$N$, VGS, a homogeneous sequential chain, a cross-model chain, MAgICoRe \citep{chen2025magicore}, multi-agent debate \citep{du2024debate}, Mixture-of-Agents (MoA) \citep{wang2025moa}, and a Hierarchical Planner-Executor (HPE; adapted from \citet{wang2023plansolve} and \citet{hong2023metagpt}). Same-model protocols use Claude Sonnet~4; cross-model variants add GPT-4o and Gemini~2.5 Flash via OpenRouter.

\paragraph{VGS (Verifier-Guided Search).}
VGS maintains a population of $N{=}4$ candidate solutions and iterates for up to $K$ evaluator calls. Each generation: (1)~score all candidates with the development evaluator; (2)~select the top-2 by score; (3)~prompt the LLM with the selected candidates and their scores (OPRO-style \citep{yang2024opro}) to produce $N$ new candidates via mutation; (4)~replace the population. The loop terminates when the evaluator-call cap is reached. VGS is a single-agent protocol, one model generates all candidates, and represents aggressive exploitation of the visible evaluator.

\paragraph{Key design note.}
Best-of-$N$ makes $N{=}8$ independent calls from one model and selects by dev score. It consumes comparable tokens to debate or MAgICoRe but with zero interaction and roughly $3{\times}$ more evaluator calls (${\sim}8$ vs.\ ${\sim}3$; Appendix~\ref{app:budget}). Despite this evaluator-call advantage, Best-of-$N$ ranks mid-pack ($-0.051$), roughly tied with cross-chain and homo-chain, while Self-Refine outranks it with only iterative single-agent refinement. Eight calls to the same model yield diminishing returns; the marginal value of a new sample depends on how different it is from existing samples, not on how many evaluator calls produced it.

\paragraph{Diversity--synthesis ablation.}
To isolate the source of MoA's relative advantage, we run a balanced $2{\times}2$ factorial that varies model-family diversity and aggregation method, with all four cells generating exactly $N{=}3$ proposals: MoA (diverse + synthesis), MoA-same-model (same + synthesis), MoA-no-synth (diverse + selection only), and Best-of-3 (same + selection only). \emph{Synthesis} means a designated model reads all three proposals and produces a combined output; \emph{selection} means choosing the highest-scoring proposal. The primary ablation uses Claude Sonnet~4 as synthesizer; a robustness check substitutes GPT-4o and Gemini~2.5 Flash (Appendix~\ref{app:multisynth}). All cells run 10~seeds on the three tasks with sufficient hidden-score variation ($N{=}120$ primary observations). The remaining tasks show too little variation to resolve differences.

\paragraph{Prescriptive protocol (derived from Section~\ref{sec:finding2}).}
Best-of-5-Diverse generates one proposal from each of five model families (Claude, GPT-4o, Gemini~2.5 Flash, DeepSeek-V3, Qwen-2.5-72B) independently, with no synthesis or interaction, and selects the highest-scoring proposal by dev score. It is derived from the diversity tax analysis: if interaction erodes diversity, a natural protocol should maximize model-family diversity while eliminating interaction entirely.

\paragraph{Model-family-diversity ablation.}
Protocols mixing model families for MAgICoRe, Debate, and HPE run in both frontier (Claude$+$GPT-4o$+$Gemini) and open-weight (DeepSeek-V3$+$Llama-3.3-70B$+$Qwen-2.5-72B) configurations, giving twelve conditions in a $2{\times}2$ (frontier vs.\ OSS) $\times$ (homogeneous vs.\ mixed) design. OSS MEG is computed against DeepSeek-specific single-agent controls to keep the denominator family-matched.

% ══════════════════════════════════════════════════════════════════════════════
\section{Results}
\label{sec:results}

\subsection{Adding agents does not help}
\label{sec:finding1}

\begin{figure}[t]
\centering
\includegraphics[width=\linewidth]{figures/agg_meg_core_bar.pdf}
\caption{Aggregate MEG with bootstrap 95\% CIs (9 tasks, hidden evaluator). No core Claude-family protocol achieves positive MEG. Diverse MoA (the only protocol using three different model families) is the only one whose CI spans zero. Agent count does not explain the ranking ($\rho_s{\approx}0$); information diversity does.}
\label{fig:meg}
\end{figure}

Figure~\ref{fig:meg} shows the aggregate MEG ranking across all nine tasks:

\begin{itemize}
\item \textbf{MoA-diverse ($-0.022$)} is the only protocol consistent with null, and it is the only one using three different model families.
\item \textbf{Self-Refine ($-0.038$)} is next; its CI narrowly excludes zero, making it weakly negative.
\item \textbf{Best-of-$N$ ($-0.051$)}, \textbf{cross-chain ($-0.053$)}, and \textbf{homo-chain ($-0.055$)} are roughly tied mid-pack. Best-of-$N$ uses ${\sim}8$ evaluator calls vs.\ ${\sim}3$ for most multi-agent protocols (Section~\ref{sec:protocols}).
\item \textbf{MoA-same-model ($-0.091$)} is clearly negative---same structure and agent count as diverse MoA, but all proposers are Claude.
\item \textbf{HPE ($-0.225$)} is worst: the planner's subtask decompositions are frequently not easily recombined into a single schema-valid artifact.
\end{itemize}

No core Claude-family protocol achieves a positive MEG estimate (per-task breakdown in Appendix~\ref{app:megtable}). A Spearman rank correlation of AggMEG on agent count across the 10 protocols yields $\rho_s{\approx}0$ ($p{>}0.9$): the ranking is not explained by agent count. Self-Refine (1 agent) outranks debate, MAgICoRe, and HPE (all multi-agent), while Best-of-$N$ (1 agent, 8 calls) is roughly tied with cross-chain (2 agents, different models).

\subsection{Interaction destroys the diversity it needs}
\label{sec:finding2}

\begin{table}[t]
\caption{Balanced $2{\times}2$ diversity--synthesis ablation ($Q$-normalized hidden scores, three tasks, all cells $N{=}3$ proposals, 10 seeds each, $N{=}120$ total observations). Model-family diversity (columns) is the only significant factor; the aggregation mechanism (rows) is not.}
\label{tab:ablation22}
\centering
\small
\begin{tabular}{l l c c}
\toprule
& & Same model family & Diverse model families \\
\midrule
\multicolumn{4}{l}{\textit{Difference Bases}} \\
 \quad & Synthesis (combined output) & $0.07$ & $\mathbf{0.29}$ \\
 \quad & Selection (best-of-proposals) & $0.20$ & $\mathbf{0.24}$ \\
\midrule
\multicolumn{4}{l}{\textit{Erd\H{o}s Overlap}} \\
\quad & Synthesis (combined output) & $0.07$ & $\mathbf{0.50}$ \\
\quad & Selection (best-of-proposals) & $0.03$ & $\mathbf{0.57}$ \\
\midrule
\multicolumn{4}{l}{\textit{Molecule QED}} \\
\quad & Synthesis (combined output) & $0.38$ & $\mathbf{0.44}$ \\
\quad & Selection (best-of-proposals) & $0.31$ & $0.31$ \\
\bottomrule
\end{tabular}
\end{table}

Table~\ref{tab:ablation22} compares model-family diversity and synthesis on the three tasks with sufficient hidden-score variation:

\begin{itemize}
\item \textbf{Model-family diversity is the larger effect.} On DiffBases and Erd\H{o}s, diverse models outscore same-model in every row ($0.50$ vs.\ $0.07$ with synthesis on Erd\H{o}s; $0.29$ vs.\ $0.07$ on DiffBases). The diversity effect ($+0.43$ on Erd\H{o}s) dwarfs the synthesis effect on these tasks.
\item \textbf{Synthesis is not uniformly helpful.} On Erd\H{o}s, selection outperforms synthesis for diverse models ($0.57 > 0.50$), suggesting the synthesizer partially averages out complementary proposals. On DiffBases, synthesis helps ($0.29 > 0.24$). The direction depends on whether proposals are compositionally compatible.
\item \textbf{Task composability matters.} On MolQED, diversity without synthesis produces the same mean as Best-of-3 ($0.31$), because the diverse models generate chemically incompatible structures. The diversity advantage appears only when the task supports composition of independent contributions (Section~\ref{sec:discussion}).
\end{itemize}

A linear model with task fixed effects ($Q \sim \text{diversity} \times \text{synthesis} + \text{task}$; $N{=}120$, $R^{2}{=}0.21$) confirms this quantitatively. The diversity coefficient is $+0.19$ (CI $[+0.06, +0.31]$); the synthesis coefficient is $-0.01$ (CI $[-0.11, +0.09]$); the interaction is not significant (CI $[-0.13, +0.22]$). The $R^{2}$ is modest because within-cell stochasticity (LLM sampling variance across seeds) dominates; the factorial design targets a specific hypothesis about main effects, not total variance explained, and the diversity coefficient is significant despite accounting for only part of the total variation. In 99.9\% of 10{,}000 bootstrap resamples, diversity exceeds synthesis ($\Delta{=}+0.20$, CI $[+0.07, +0.32]$). A random-intercept REML model with task as the grouping factor gives the same fixed-effect estimates, confirming that the result is not an artifact of task dummies (Appendix~\ref{app:robustness}). All four cells use $N{=}3$ proposals, eliminating the team-size confound present in the unmatched Best-of-8 comparison. A multi-synthesizer check replaces the Claude aggregator with GPT-4o and Gemini~2.5 Flash; the diversity advantage on DiffBases is $+0.22$ (Claude), $+0.14$ (GPT-4o), and $+0.17$ (Gemini), confirming that the effect is not an artifact of the synthesizer model (Appendix~\ref{app:multisynth}).

\textbf{Sensitivity to Erd\H{o}s.} Erd\H{o}s contributes the largest per-task diversity effect ($+0.48$ vs.\ $+0.13$ for DiffBases and $+0.03$ for MolQED). Excluding Erd\H{o}s, the aggregate diversity coefficient drops to $+0.076$ (CI $[-0.014, +0.168]$), no longer significant at the 5\% level but still positive and trending. On DiffBases alone, the effect is $+0.13$ (CI $[-0.02, +0.28]$). The finding is therefore driven primarily by Erd\H{o}s; DiffBases provides directional support but is individually underpowered at $n{=}10$ seeds. This is a genuine limitation of the current task suite: the diversity claim rests on two tasks, one dominant. The per-task view is more informative than the aggregate for this reason. Budget consumption across the four $2{\times}2$ cells is comparable (Diverse+Synth: 11.3k tokens, Same+Synth: 7.4k, Diverse+Select: 10.3k, Same+Select: 4.8k; all well below the 200k cap).

\paragraph{The diversity tax, measured directly.}
If interaction erodes diversity, we should see it in MIG (Eq.~\ref{eq:mig}): the score change from adding the interaction step, holding the agent team constant. Across five protocol families on eight tasks (Appendix~\ref{app:mig}), the pattern is stark. Every mixed-model protocol that makes agents communicate pays the tax: MAgICoRe ($-0.035$), Debate ($-0.078$), HPE ($-0.145$). The more structured the interaction, the heavier the cost. Only MoA has positive mixed MIG ($+0.016$), because it is composition without interaction: the synthesizer reads three independent proposals, and no proposer ever sees another's work. With same-model agents, where there is less diversity to lose, interaction adds modest value (MAgICoRe $+0.044$, Homo-chain $+0.051$), but this does not translate to positive MEG because coordination overhead consumes the budget advantage. \emph{Note}: homo and mixed MIG use different baselines (single-shot vs.\ best-of-three families), so cross-column magnitudes are not directly comparable; the within-column sign patterns are the primary evidence (Limitation~ix).

Diversity helps on two of the three tasks and is absent on the third, suggesting that the diversity advantage requires independent partial solutions to be composable. MoA-diverse is near null in aggregate because its advantage on DiffBases and Erd\H{o}s is offset by MolQED and the low-spread tasks where diversity cannot be exploited. A pre-specified composability classification (Appendix~\ref{app:composability}), written from task structure alone before examining these results, correctly predicts the pattern for 8 of 9 tasks but misclassifies MolQED as HIGH; modular domain structure does not always map onto the model's output representation.

\paragraph{What the $2{\times}2$ and MIG together show.}
The factorial isolates the source of value: independence across model families. The MIG decomposition identifies the mechanism that destroys it: interaction. Together, they reveal a cost at the heart of multi-agent protocol design. Debate, MAgICoRe, and HPE improve within-model quality (positive MIG for same-model agents) but destroy cross-model complementarity (negative MIG for diverse agents). The more structured the interaction, the heavier the tax. This is the ensemble-diversity phenomenon in a new setting: the value of combining multiple sources lies in their independence, and any process that creates correlation---social influence in human groups \citep{lorenz2011social}, iterative debate among LLMs---erodes it. Multi-agent debate is social influence for language models.

Three independent evidence streams converge on this conclusion: (1)~the $2{\times}2$ factorial identifies diversity as the only significant factor across $N{=}120$ observations; (2)~the MIG decomposition shows that 4 of 5 mixed-model protocol families pay a negative interaction tax, spanning 8 tasks and 32 conditions with a consistent sign pattern; (3)~the 41-variant sweep across four model families finds only one protocol consistent with null (MoA-diverse), and it is the only one using diverse models without interaction. No single stream is conclusive alone---the $2{\times}2$ rests on two strong tasks, MIG uses asymmetric baselines (Limitation~ix), and the sweep is exploratory---but their agreement is the primary evidence for the diversity tax.

\paragraph{Exploratory design probe: Best-of-5-Diverse.}
The diversity tax predicts that maximizing model-family diversity while eliminating interaction should preserve the diversity advantage. We test this with a \emph{Best-of-5-Diverse} protocol: five models from different families (Claude, GPT-4o, Gemini~2.5 Flash, DeepSeek-V3, Qwen-2.5-72B) each generate one proposal independently; the highest-scoring proposal is selected with no synthesis step. On the three ablation tasks (10 seeds each, controls extended to $n{=}10$ seeds for tighter baselines), Best-of-5-Diverse is the only protocol with positive MEG on all three: Erd\H{o}s $+0.33$ (CI $[+0.15, +0.51]$), DiffBases $+0.11$ (CI spanning zero), MolQED $+0.05$ (CI spanning zero); aggregate MEG $= +0.16$ (CI $[+0.05, +0.25]$). No other protocol in the benchmark---across 41 variants and four model families---achieves positive MEG on all three ablation tasks. The aggregate CI excludes zero, consistent with the diversity tax prediction; Erd\H{o}s is individually significant but DiffBases and MolQED are not. The protocol was designed after the $2{\times}2$ analysis and should be treated as exploratory.

\subsection{Hidden evaluation as a methodological safeguard}
\label{sec:finding3}

On most tasks, dev and hidden rankings agree ($\rho_s{>}0.9$); hidden evaluation does not overturn the overall negative finding. The one clear divergence is \textbf{Molecule QED} ($\rho_s{=}0.939$, top-3 overlap $2/3$): agents that hill-climb QED on the visible evaluator produce molecules penalized by the hidden synthesizability criterion \citep{bickerton2012qed,ertl2009sascore}---the agent maximizes a proxy rather than the true objective, changing the top-3 set but not the winner. Two other tasks show minor rank shifts in near-zero regimes (MaxCut: $\rho_s{=}0.891$, all $Q{\approx}0$; DiffBases: VGS moves from last on dev to \#3 on hidden; Appendix~\ref{app:kcurve}). The evidence supports hidden evaluation as a prudent methodological safeguard---the ``held-out test set'' for optimization benchmarks---rather than a dramatic correction on this suite.

\subsection{The diversity tax generalizes across model families}
\label{sec:finding4}

\begin{figure}[t]
\centering
\includegraphics[width=\linewidth]{figures/meg_heatmap_primary.pdf}
\caption{Per-task MEG by protocol (10 core protocols, 9 tasks). The task-conditional structure is a prominent pattern: DiffBases and Erd\H{o}s show isolated positive entries; the low-spread tasks (FlatPoly, TSP-100, LJ) show uniform near-zero MEG; TSP-50 shows modest discrimination (VGS and Best-of-$N$ lead).}
\label{fig:heatmap}
\end{figure}

Figure~\ref{fig:heatmap} shows the per-task structure. An \emph{exploratory} extension to 41 protocol variants across GPT-4o, Gemini~2.5 Flash, and DeepSeek-V3 (Appendix~\ref{app:allvariants}) finds only one positive entry: Gemini MAgICoRe ($+0.011$, CI spanning zero). Agent count does not explain performance in any model family. A model-family-diversity ablation across frontier and open-weight configurations (Appendix~\ref{app:braindiv}) echoes the diversity tax: mixing diverse OSS models reduces the performance gap for Debate ($+0.029$) and HPE ($+0.187$) but not MAgICoRe ($-0.013$); no condition crosses the single-agent baseline.

% ══════════════════════════════════════════════════════════════════════════════
\section{Discussion}
\label{sec:discussion}

\paragraph{The diversity tax: structural analogy.}
The ensemble-learning ambiguity decomposition \citep{krogh1994neural} shows that ensemble accuracy equals average individual accuracy plus a diversity term; adding a correlated member improves average quality but shrinks the diversity term, and the net can be negative. The formal decomposition assumes squared loss and weighted averaging, neither of which holds exactly here; the analogy is \emph{structural} rather than mathematical---interaction creates correlation, and correlation erodes the value of combination. The gradient across protocol families (Section~\ref{sec:finding2}) mirrors this structure: HPE ($-0.145$) imposes the most coordination and pays the heaviest tax; MoA ($+0.016$) avoids it entirely by preserving independence.

\paragraph{Composition is not interaction.}
The multi-agent literature has largely conflated two distinct operations: \emph{interaction} (iterative exchange between agents---debate rounds, refinement loops, hierarchical instruction) and \emph{composition} (combining independent contributions in one shot). MoA's one-shot synthesis is the only mixed-model protocol with positive MIG ($+0.016$) because it is composition without interaction: the synthesizer reads three independent proposals and produces a combined artifact, but no proposer ever sees another proposer's work. Independence is preserved by construction. The implication for protocol design is that the bottleneck is not generating diverse proposals---different model families already provide that---but composing them without introducing convergence pressure. Protocols that merge partial solutions without iterative back-and-forth are more likely to preserve the diversity advantage.

\paragraph{Solution-space diversity vs.\ score diversity.}
A natural question is whether model-family diversity produces measurably different \emph{solutions}, not just different \emph{scores}. We compute pairwise solution distances within each $2{\times}2$ cell---Jaccard distance for DiffBases (integer sets), cosine distance for Erd\H{o}s (value vectors), and normalized Levenshtein distance for MolQED (SMILES strings)---and compare diverse-model vs.\ same-model arms (bootstrap 95\% CIs, 10{,}000 resamples).

The three tasks tell complementary stories. On \textbf{DiffBases}, diverse models produce marginally more dissimilar solutions ($\Delta_{\mathrm{sol}}{=}{+}0.035$, CI $[-0.018, +0.085]$, n.s.), while the score diversity is large ($\Delta_Q{=}{+}0.13$). On \textbf{Erd\H{o}s}, diverse models produce \emph{less} cosine-diverse solutions ($\Delta_{\mathrm{sol}}{=}{-}0.099$, CI $[-0.180, -0.019]$), despite the largest score diversity of any task ($\Delta_Q{=}{+}0.48$). On \textbf{MolQED}, diverse models produce significantly \emph{more} diverse SMILES strings ($\Delta_{\mathrm{sol}}{=}{+}0.129$, CI $[+0.063, +0.195]$), yet the score diversity is null ($\Delta_Q{=}{+}0.03$, n.s.). The pattern aligns with the composability hypothesis: on composable tasks, the diversity advantage comes from complementary \emph{quality} regions rather than geometric spread in solution space; on the non-composable task, even large solution-space diversity does not translate to score improvement because diverse solutions cannot be usefully composed. The ensemble interpretation holds: the value of independence lies in reducing correlated errors, not in maximizing output distance.

\paragraph{When does diversity help? A composability hypothesis.}
The diversity advantage is not universal: it appears on DiffBases and Erd\H{o}s (where different models find structurally distinct valid components that can be merged) but not on MolQED (where diverse models generate chemically incompatible molecules). A pre-specified composability classification (Appendix~\ref{app:composability}) correctly predicts which tasks benefit for 8 of 9; the MolQED miss reveals that the criterion must account for solution-representation compatibility, not just task-level modularity. We hypothesize that \emph{model-family diversity improves multi-agent performance if and only if the task admits composition of independent partial solutions}. The hypothesis is consistent with the current data (2/3 ablation tasks, 8/9 full suite, directional support from three out-of-suite tasks) but has not been tested on a held-out task set; confirming or refuting it requires no access to our benchmark.

\paragraph{MEG vs.\ MIG: different questions.}
MEG asks whether a protocol beats the best single-agent control (a practitioner question); MIG asks whether the interaction step adds value over the best individual agent in the same team (a mechanistic question). MoA's positive MIG ($+0.016$) but negative MEG ($-0.022$) illustrates the gap: one-shot synthesis extracts value from diversity, but coordination overhead leaves the full protocol below a single-agent control that spends its budget differently. The MIG baselines are asymmetric across configurations (Limitation~ix), so within-configuration comparisons are more reliable than cross-configuration ones.

\paragraph{Hidden evaluation as a guard against proxy overfit.}
On at least one task (MolQED), protocols that hill-climb the visible scorer are penalized by the hidden criterion. This is the same failure mode as reward hacking in RLHF: the agent maximizes a proxy rather than the true objective. An out-of-suite probe (Appendix~\ref{app:external}) provides suggestive---though not conclusive---corroboration: synthesis shows a large significant effect on Partition under visible-only scoring ($+857$, CI $[+341, +1364]$) but is non-significant under hidden evaluation on the main benchmark. Because these are different tasks, the comparison cannot isolate hidden evaluation as the cause; the pattern is consistent with hidden evaluation suppressing scorer-specific gains but could also reflect task-level differences. Hidden evaluation is the ``held-out test set'' for optimization benchmarks, and both the MolQED ranking inversion and the suggestive synthesis contrast motivate it as standard practice.

\paragraph{Budget scaling does not rescue the negative finding.}
A $2{\times}$ budget experiment on three discriminative tasks (Appendix~\ref{app:budget-scaling}), including MoA-diverse, finds that single-agent controls scale faster than all protocols tested: doubling budget does not rescue a positive MEG. A positive MEG regime may exist at much higher budgets or with protocols designed for inter-agent information exchange rather than adapted from reasoning-chain methods, but the current evidence points away from it.

\paragraph{Limitations.}
(i)~Of nine tasks, four have sufficient hidden-score variation to separate protocols (DiffBases, Erd\H{o}s, MolQED, TSP-50); the diversity finding from the $2{\times}2$ is strongest on Erd\H{o}s ($+0.48$), moderate on DiffBases ($+0.13$), and absent on MolQED. Excluding Erd\H{o}s, the aggregate diversity coefficient drops to $+0.076$ (CI spanning zero); the per-task view is therefore more informative than the aggregate (Section~\ref{sec:finding2}). Aggregate rankings are sensitive to task composition: under leave-one-task-out, MoA's rank ranges from 1 to 2; under median aggregation, Best-of-$N$ ranks first instead of MoA. MoA and Best-of-$N$ are the only protocols whose jackknife $\pm2$SE intervals span zero (Appendix~\ref{app:robustness}). The per-task view is primary, and the claim is a benchmark-specific characterization, not a general law. (ii)~The suite remains heavier on combinatorial tasks than on open-ended scientific domains. An external probe on three out-of-suite tasks (Appendix~\ref{app:external}) provides directionally consistent but non-significant ($n{=}5$) evidence for the composability prediction; the diversity finding rests primarily on the original benchmark's $2{\times}2$ factorial. (iii)~Several tasks show little score variation ($Q \approx 0$) under the current budget and remain unresolved at $2{\times}$ (Appendix~\ref{app:budget-scaling}). (iv)~Best-of-$N$ uses ${\sim}3{\times}$ more evaluator calls than most multi-agent protocols (${\sim}8$ vs.\ ${\sim}3$ calls), yet ranks mid-pack (tied with cross-chain); the comparison does not cleanly separate evaluator-call intensity from protocol structure. (v)~HPE open-weight data is sparse ($<5$ seeds on some tasks). (vi)~We use the term \emph{structured optimization} rather than \emph{scientific reasoning} throughout, because scores are determined by exact objectives applied to artifacts, not by assessment of reasoning. Several tasks (MaxCut, TSP, Flat Polynomials) are standard mathematical optimization benchmarks; LJ clusters and Molecule QED are the most domain-specific. (vii)~All tasks, evaluators, protocols, and analysis originate from a single author; independent replication would strengthen the claims. (viii)~Multi-agent protocols are our implementations, not the original authors' code; suboptimal prompt engineering could understate the best-case performance of debate, MAgICoRe, or HPE. Our Self-Refine variant replaces the self-critique loop of \citet{madaan2023selfrefine} with direct verifier feedback, making it a stronger control baseline; the negative MEG finding for interaction-heavy protocols is therefore conservative with respect to this deviation. (ix)~Extended MIG uses single-shot as the no-interact baseline for homo protocols, but best-of-three model families for mixed protocols; the asymmetric baseline means homo and mixed MIG are not directly comparable in magnitude (Appendix~\ref{app:mig}). (x)~Bootstrap CIs are not adjusted for multiple comparisons across 41 variants $\times$ 9 tasks; the main claims use aggregate MEG (one CI per protocol), but per-task positive entries in the appendix should be interpreted cautiously.

% ══════════════════════════════════════════════════════════════════════════════
\section{Conclusion}
\label{sec:conclusion}

The thing that helps multi-agent systems is not adding agents, refining answers, or designing better debates. It is preserving independence. Under matched budget and hidden evaluation, model-family diversity is the only significant factor ($+0.19$, CI $[+0.06, +0.31]$), and every interaction mechanism we tested destroys it. The benchmark reveals a consistent diversity tax: debate, refinement, and planning push agents toward consensus, erasing the complementary information that makes diverse configurations valuable. Only one-shot composition avoids the tax, because no agent ever sees another's work.

Three implications follow. \textbf{For practitioners}: on the discriminative tasks, three calls to three different models beat eight calls to one. An exploratory extension to five model families suggests this advantage scales with the number of independent families (Section~\ref{sec:finding2}). The marginal value of a new solution depends on how different it is from what you already have, not on how much compute produced it. \textbf{For benchmark designers}: hidden evaluation is the held-out test set for optimization benchmarks. Without it, visible-evaluator convergence masquerades as collaboration. \textbf{For protocol designers}: composition and interaction are not the same thing. The field has focused on how agents should talk to each other. The evidence here suggests the bottleneck is the opposite---how to combine independent contributions \emph{without} making them talk.

% ══════════════════════════════════════════════════════════════════════════════
\section*{Reproducibility Statement}
All experiments use a frozen testbed snapshot (\textsc{Veil} v1.0) with pinned model identifiers, deterministic verifiers, fixed normalization anchors, and logged budget traces. Task instances, development and hidden evaluators, prompt templates, and protocol configurations are version-controlled. Seeds are explicit integers (1--10 for MoA and $2{\times}2$ arms; 1--10 for single-agent controls on the three ablation tasks; 1--5 elsewhere); stochasticity enters only through LLM sampling temperature (fixed at $T{=}0.7$ for all models) and provider-side nondeterminism. The hidden evaluator is deterministic given a submission artifact. Budget caps ($T{=}200\text{K}$ tokens, $W{=}600\text{s}$, $C{=}30\text{s}$ CPU, $K{=}25$ calls) are enforced per cell and logged. One mid-run model substitution, Gemini preview to GA release, is documented in Section~\ref{sec:task-suite}. The release artifact includes: (1)~all frozen task instances with development and hidden evaluators packaged as deterministic Python functions; (2)~the protocol harness with budget enforcement, submission, and MEG computation; (3)~prompt templates for all 10 core protocols and 41 variants; (4)~analysis scripts that reproduce every table and figure from the result cache; (5)~the complete result cache (per-cell JSON with raw and hidden scores, budget traces, and submission artifacts) enabling provider-reproducibility checks without re-running API calls. Hidden evaluators are included but architecturally isolated: they are never imported or callable by protocol code during search. An anonymized repository is available at submission.

\section*{Ethics Statement}
This work evaluates multi-agent LLM protocols on structured optimization tasks using deterministic programmatic verifiers. No human subjects, private data, or personally identifiable information are involved. All tasks are drawn from published mathematical and structured optimization problems with no dual-use concerns. The primary risk this work addresses is methodological: benchmarks relying solely on visible-evaluator scoring can overstate multi-agent gains, leading practitioners to deploy unnecessarily complex and costly architectures.

% ══════════════════════════════════════════════════════════════════════════════
\appendix

\section{Verifier Integrity Audit}
\label{app:verifier-audit}

Prior to the experiment launch and again post-collection, we audited all development and hidden verifier implementations for robustness against malformed inputs. Three issues were found in two tasks; all were fixed before the MEG analysis was run.

\paragraph{V1: Circle Packing (both verifiers).}
The original verifier used a bare \texttt{assert} to check input shape: \texttt{assert circles.shape == (20,\,3)}. Malformed submissions raised an unhandled \texttt{AssertionError} inside the evaluation subprocess; the platform exception handler mapped this to a score of zero, but exception-based sentinels are fragile under Python version drift and sandbox policy changes. Fixed by replacing the assert with an explicit \texttt{try}/\texttt{except} and shape check returning \texttt{float('-inf')}.

\paragraph{V2: Lennard-Jones $n=41$ (both verifiers).}
Same pattern: \texttt{assert pos.shape == (41,\,3)} raised \texttt{AssertionError} on malformed position arrays. Fixed identically, returning the task sentinel \texttt{1e30} (large positive, consistent with the minimize direction and distinguished from any physically valid energy).

\paragraph{V3: Lennard-Jones $n=41$ energy validity gate.}
The original verifier did not check the sign of the computed total LJ energy. A non-negative total energy indicates atoms at overlap distances where the $r^{-12}$ repulsion term dominates, producing values in $[10^{60}, 10^{180}]$ that corrupt the $Q$ normalization scale if treated as valid. Fixed by returning \texttt{1e30} whenever the computed energy is non-negative. The gate is applied to both the bare LJ energy and the Axilrod-Teller correction in the hidden verifier.

\paragraph{Sensitivity.}
All three fixes affect only malformed or physically invalid submissions. Every submission in the original analysis dataset is well-formed and passes the original assert gates; no LJ submission has a non-negative energy. No score in the MEG analysis table changes as a result of the audit.

\section{Score Normalization Anchors}
\label{app:normalization}

Per-task $Q$ scores are computed using two frozen anchors: $s_\mathrm{bad}$ (a mediocre reference, $Q{=}0$) and $s_\mathrm{ref}$ (a target good-performance level, $Q{=}1$). For minimize tasks: $Q(s) = (s_\mathrm{bad} - s)/(s_\mathrm{bad} - s_\mathrm{ref})$; for maximize tasks: $Q(s) = (s - s_\mathrm{bad})/(s_\mathrm{ref} - s_\mathrm{bad})$; both clamped to $[0,1]$. Scores worse than $s_\mathrm{bad}$ and invalid/timed-out submissions receive $Q{=}0$. Anchors (v1.0, frozen before sweep):

\begin{center}
\small
\begin{tabular}{lrrl}
\toprule
Task & $s_\mathrm{bad}$ ($Q{=}0$) & $s_\mathrm{ref}$ ($Q{=}1$) & Direction \\
\midrule
MaxCut $G(200,0.3)$      & 3139     & 3517        & maximize \\
Circle Packing $n{=}20$  & 0.5039   & 2.0         & maximize \\
Difference Bases         & 14.0     & 1.0         & minimize \\
Flat Polynomials deg.~50 & 2.0489   & 1.0         & minimize \\
TSP-100                  & 50835    & 7517        & minimize \\
Lennard-Jones $n{=}41$   & $-136.0$ & $-198.061$  & minimize \\
Erd\H{o}s overlap        & 0.2939   & 0.2321      & minimize \\
Molecule QED             & 0.60     & 1.0         & maximize \\
TSP-50                   & 25400    & 5500        & minimize \\
\bottomrule
\end{tabular}
\end{center}

All observed core-protocol scores on MaxCut fall below $s_\mathrm{bad}{=}3139$, so $Q{=}0$ for every protocol on that task; MaxCut is therefore diagnostic only via raw-rank comparison (Section~\ref{sec:finding3}), not via $Q$-based MEG.

\section{Model-Family-Diversity Ablation: Full Results}
\label{app:braindiv}

\begin{table}[h]
\caption{Model-family-diversity ablation: aggregate MEG (mean over 8--9 tasks, 5 seeds per cell; HPE OSS sparse). OSS MEG computed against DeepSeek-specific single-agent controls. All 12 cells are negative. $\Delta_\mathrm{mix,OSS}$ (bold where positive) shows mixing diverse OSS models reduces the gap for Debate and HPE but not MAgICoRe. $\Delta_\mathrm{mix}$: mixed$-$homo (frontier). $\Delta_\mathrm{cap}$: homo-OSS$-$homo-frontier. $\Delta_\mathrm{mix,OSS}$: mixed-OSS$-$homo-OSS.}
\label{tab:braindiv}
\centering
\small
\begin{tabular}{l r r r r r r r}
\toprule
Protocol & \multicolumn{2}{c}{Frontier} & \multicolumn{2}{c}{Open-Source} & $\Delta_\mathrm{mix}$ & $\Delta_\mathrm{cap}$ & $\Delta_\mathrm{mix,OSS}$ \\
\cmidrule(lr){2-3}\cmidrule(lr){4-5}
 & Homo & Mixed & Homo & Mixed & & & \\
\midrule
MAgICoRe & $-0.063$ & $-0.068$ & $-0.044$ & $-0.058$ & $-0.005$ & $+0.019$ & $-0.013$ \\
Debate   & $-0.088$ & $-0.111$ & $-0.137$ & $-0.108$ & $-0.023$ & $-0.049$ & $\mathbf{+0.029}$ \\
HPE      & $-0.225$ & $-0.176$ & $-0.438^*$ & $-0.250^*$ & $+0.049$ & $-0.213$ & $\mathbf{+0.187}^*$ \\
\bottomrule
\end{tabular}

\smallskip\noindent\small $^*$Sparse: HPE OSS cells have $n{<}5$ seeds; treat as indicative only.
\end{table}

No cell is positive in either the frontier or OSS arm. $\Delta_\mathrm{mix,OSS}$ is positive for Debate ($+0.029$) and HPE ($+0.187$), suggesting mixing diverse OSS model families can reduce the performance gap relative to homogeneous DeepSeek for those families. MAgICoRe is a counterexample ($-0.013$). Within the frontier arm, mixing provides no consistent benefit, suggesting frontier models may already share enough training-distribution overlap to limit the diversity gain.

\section{Budget Consumption}
\label{app:budget}

All protocols run under the same hard caps per cell: $T{=}200{,}000$ tokens, $W{=}600\,\text{s}$ wall-clock, $C{=}30\,\text{s}$ evaluator CPU, $K{=}25$ evaluator calls. \textbf{Budget matching is on caps, not on realized usage}: single-shot uses 1.3k tokens while VGS uses 24.4k. Protocols that use fewer calls or tokens do so because their design does not require them, not because they are constrained. No protocol exhausted any cap component on average.

\begin{table}[h]
\caption{Budget consumption per protocol (mean across 9 tasks; 5 seeds per cell for most protocols, 10 for MoA arms).}
\label{tab:budget}
\centering
\small
\begin{tabular}{l r r r r}
\toprule
Protocol & Tokens (k) & Wall (s) & Eval CPU (s) & Eval Calls \\
\midrule
Single-shot          &      1.3 &    11.4 &         0.47 &        1.0 \\
Self-Refine          &      6.4 &    41.1 &         1.35 &        3.9 \\
Best-of-$N$          &     10.3 &    91.3 &         2.67 &        7.9 \\
VGS                  &     24.4 &   106.9 &         3.92 &       10.4 \\
Homo-chain           &      5.6 &    33.0 &         1.39 &        4.0 \\
Cross-chain          &      6.0 &    26.3 &         1.08 &        2.9 \\
MAgICoRe             &     12.7 &    84.2 &         1.08 &        2.8 \\
Debate               &     13.4 &    98.6 &         1.09 &        2.9 \\
MoA                  &      8.1 &    42.2 &         1.45 &        3.9 \\
HPE                  &     21.4 &   160.0 &         0.70 &        1.9 \\
\midrule
Cap (per cell)        & 200.0 &   600.0 &        30.00 &       25.0 \\
\bottomrule
\end{tabular}
\end{table}

\section{Per-Task MEG Table}
\label{app:megtable}

\begin{table}[h]
\caption{Per-task MEG summary for the 10 core protocols (9 tasks, hidden evaluator; MoA uses seeds 1--10, all other protocols seeds 1--5). Bold = positive. Aggregate is the unweighted mean across all nine tasks.}
\centering
\small
\begin{tabular}{l r r r r r r r}
\toprule
Protocol & DiffBases & Erd\H{o}s & MolQED & TSP-50 & CircPack & Low-var$^\dagger$ & AggMEG \\
\midrule
MoA$^{n=10}$ & \textbf{+0.090} & \textbf{+0.021} & $-0.013$ & $-0.054$ & $-0.238$ & $-0.001$ & $-0.022$ \\
Self-Refine$^*$ & $-0.110$ & $0.000$ & $-0.049$ & $-0.077$ & $-0.027$ & $-0.019$ & $-0.038$ \\
Best-of-$N$$^*$ & $0.000$ & $-0.371$ & $0.000$ & $-0.003$ & $0.000$ & $-0.021$ & $-0.051$ \\
Cross-chain & $-0.113$ & $-0.134$ & $-0.102$ & $-0.021$ & $-0.027$ & $-0.021$ & $-0.053$ \\
Homo-chain & $-0.197$ & $-0.048$ & $-0.122$ & $-0.038$ & $0.000$ & $-0.024$ & $-0.055$ \\
MAgICoRe & $-0.144$ & $-0.149$ & $-0.137$ & $-0.046$ & $0.000$ & $-0.022$ & $-0.063$ \\
VGS$^*$ & $-0.131$ & $-0.476$ & $-0.041$ & $0.000$ & $0.000$ & $0.000$ & $-0.072$ \\
Debate & \textbf{+0.112} & $-0.334$ & \textbf{+0.024} & $-0.018$ & $-0.522$ & $-0.014$ & $-0.088$ \\
Single-shot & $-0.197$ & $-0.476$ & $-0.110$ & $-0.090$ & $0.000$ & $-0.022$ & $-0.107$ \\
HPE$^*$ & $-0.197^*$ & $-0.476^{**}$ & $-0.240$ & $-0.015$ & $-1.000$ & $-0.024$ & $-0.225$ \\
\bottomrule
\end{tabular}
\vspace{2pt}
{\footnotesize $^{n=10}$MoA row uses 10 seeds (seeds 1--10); all other rows use 5 seeds. $^\dagger$Low-var: mean MEG over MaxCut (diagnostic, all $Q{=}0$), FlatPoly, TSP-100, LJ-n41 (low-spread tasks where all protocols score $Q \approx 0$). $^*$\,Some task cells have $<5$ valid seeds (Self-Refine: 1 on TSP; BoN: 3 on TSP; VGS: 4 on TSP; HPE: 4 on DiffBases). $^{**}$\,1 valid seed (HPE on Erd\H{o}s). Entries of $0.000$ for control protocols (BoN, VGS, SR) indicate that protocol was the binding control on that task (MEG $=$ itself $-$ itself $= 0$ by construction).}
\end{table}

\section{MaxCut Diagnostic}
\label{app:kcurve}

\begin{figure}[h]
\centering
\includegraphics[width=\linewidth]{figures/maxcut_rank_inversion.pdf}
\caption{MaxCut dev/hidden rank comparison (10 core protocols, raw scores). VGS ranks first on the public 80\%-edge evaluator but drops on the hidden full-graph evaluator; cross-chain rises to rank~1. All protocols score at or below a random partition ($Q{=}0$), so this is rank instability in a regime where all protocols score $Q{=}0$.}
\label{fig:maxcut}
\end{figure}

\begin{figure}[h]
\centering
\includegraphics[width=0.75\linewidth]{figures/kcurve_maxcut.pdf}
\caption{MaxCut K-curve: dev and hidden scores as a function of evaluator-call cap $K$ (VGS, 5 seeds per $K$). Both curves rise together from $K{=}1$ to $K{=}25$ ($\Delta_{\mathrm{dev}}{=}+29$, $\Delta_{\mathrm{hidden}}{=}+54$), indicating that the dev evaluator is not gameable. The hidden gain exceeds the dev gain, ruling out evaluator exploitation.}
\label{fig:kcurve}
\end{figure}

\section{Full Protocol Variant Table}
\label{app:allvariants}

\begin{table}[h]
\caption{Aggregate MEG for all 41 protocol variants, grouped by model family. The 10 core Claude protocols (top group) are the primary comparison; the remaining 31 variants are exploratory extensions that test generalization across model families. $k$: number of tasks with data. Sorted by AggMEG within each group. CI is bootstrap 95\%. $^{n=10}$MoA arms use 10 seeds; all others use 5. $^*$Some task cells have $<5$ valid seeds.}
\centering
\scriptsize
\begin{tabular}{l r r r}
\toprule
Protocol variant & $k$ & AggMEG & 95\% CI \\
\midrule
\multicolumn{4}{l}{\textit{Claude Sonnet~4: core protocols}} \\
MoA$^{n=10}$ & 9 & $-0.022$ & $[-0.059, +0.015]$ \\
Self-Refine & 9 & $-0.038$ & $[-0.071, -0.005]$ \\
Best-of-$N$ & 9 & $-0.051$ & $[-0.084, -0.013]$ \\
Cross-chain & 9 & $-0.053$ & $[-0.089, -0.016]$ \\
Homo-chain & 9 & $-0.055$ & $[-0.080, -0.030]$ \\
MAgICoRe & 9 & $-0.063$ & $[-0.098, -0.028]$ \\
VGS & 9 & $-0.072$ & $[-0.091, -0.048]$ \\
Debate & 9 & $-0.088$ & $[-0.142, -0.031]$ \\
Single-shot & 9 & $-0.107$ & $[-0.118, -0.094]$ \\
HPE$^*$ & 9 & $-0.225$ & $[-0.247, -0.201]$ \\
\midrule
\multicolumn{4}{l}{\textit{Claude: ablation arms}} \\
MoA-same-model$^{n=10}$ & 9 & $-0.091$ & $[-0.111, -0.068]$ \\
MoA-no-synth$^{n=10}$ & 8 & $-0.044$ & $[-0.082, -0.006]$ \\
ToT & 6 & $-0.013$ & $[-0.034, +0.008]$ \\
\midrule
\multicolumn{4}{l}{\textit{Claude: frontier mixed models}} \\
MAgICoRe-mixed & 8 & $-0.068$ & $[-0.100, -0.026]$ \\
Debate-mixed$^*$ & 8 & $-0.111$ & $[-0.151, -0.081]$ \\
HPE-mixed$^*$ & 7 & $-0.176$ & $[-0.226, -0.124]$ \\
\midrule
\multicolumn{4}{l}{\textit{GPT-4o (8 tasks, MEG vs.\ GPT-4o controls)}} \\
Debate-GPT4o & 8 & $-0.035$ & $[-0.072, +0.003]$ \\
BoN-GPT4o & 8 & $-0.041$ & $[-0.079, -0.004]$ \\
VGS-GPT4o & 8 & $-0.047$ & $[-0.083, -0.009]$ \\
SR-GPT4o & 8 & $-0.060$ & $[-0.101, -0.025]$ \\
Homo-chain-GPT4o & 8 & $-0.081$ & $[-0.123, -0.043]$ \\
MAgICoRe-GPT4o & 8 & $-0.084$ & $[-0.127, -0.043]$ \\
HPE-GPT4o$^*$ & 7 & $-0.174$ & $[-0.230, -0.117]$ \\
SS-GPT4o & 8 & $-0.189$ & $[-0.233, -0.147]$ \\
\midrule
\multicolumn{4}{l}{\textit{DeepSeek-V3 (4--5 tasks; MEG vs.\ DeepSeek controls)}} \\
BoN-DeepSeek & 5 & $-0.003$ & $[-0.063, +0.057]$ \\
VGS-DeepSeek & 5 & $-0.041$ & $[-0.079, +0.019]$ \\
MAgICoRe-DeepSeek & 5 & $-0.044$ & $[-0.109, +0.025]$ \\
SR-DeepSeek & 5 & $-0.077$ & $[-0.123, -0.016]$ \\
Debate-DeepSeek$^*$ & 5 & $-0.137$ & $[-0.208, -0.072]$ \\
HPE-DeepSeek$^*$ & 4 & $-0.438$ & $[-0.506, -0.344]$ \\
\midrule
\multicolumn{4}{l}{\textit{Gemini~2.5 Flash (8 tasks, MEG vs.\ Gemini controls)}} \\
MAgICoRe-Gemini & 8 & $+0.011$ & $[-0.060, +0.084]$ \\
VGS-Gemini & 8 & $-0.018$ & $[-0.083, +0.047]$ \\
HPE-Gemini$^*$ & 6 & $-0.038$ & $[-0.110, +0.041]$ \\
BoN-Gemini & 8 & $-0.041$ & $[-0.101, +0.025]$ \\
SR-Gemini & 8 & $-0.109$ & $[-0.166, -0.052]$ \\
Debate-Gemini & 8 & $-0.111$ & $[-0.162, -0.054]$ \\
Homo-chain-Gemini & 8 & $-0.118$ & $[-0.170, -0.055]$ \\
SS-Gemini & 8 & $-0.186$ & $[-0.216, -0.146]$ \\
\midrule
\multicolumn{4}{l}{\textit{OSS mixed models (DeepSeek+Llama+Qwen; MEG vs.\ DeepSeek controls)}} \\
MAgICoRe-mixed-OSS$^*$ & 8 & $-0.058$ & $[-0.085, -0.028]$ \\
Debate-mixed-OSS$^*$ & 8 & $-0.108$ & $[-0.146, -0.051]$ \\
HPE-mixed-OSS$^*$ & 8 & $-0.250$ & $[-0.291, -0.206]$ \\
\bottomrule
\end{tabular}
\end{table}

\section{MEG Denominator Bias Correction}
\label{app:megbias}

The MEG denominator $\max(\bar{Q}_{\mathrm{SR}}, \bar{Q}_{\mathrm{BoN}}, \bar{Q}_{\mathrm{VGS}})$ takes the maximum of three sample means, each from $n{=}5$ seeds. Because $\mathbb{E}[\max(\bar{X}_1, \bar{X}_2, \bar{X}_3)] \geq \max(\mathbb{E}[\bar{X}_1], \mathbb{E}[\bar{X}_2], \mathbb{E}[\bar{X}_3])$, the denominator is biased upward, pushing MEG conservatively negative. We estimate the bias via bootstrap: for each task, we resample the control seeds 10{,}000 times with replacement, compute the expected max-of-means, and subtract the raw max-of-means.

\begin{table}[h]
\caption{Per-task denominator bias (bootstrap, 10K resamples). Bias is positive when control means are noisy and close together; zero when controls collapse to $Q{=}0$ (MaxCut, LJ). Mean bias across nine tasks: $+0.006$.}
\label{tab:megbias}
\centering
\small
\begin{tabular}{l r r r}
\toprule
Task & Raw denom & $\mathbb{E}[\text{boot denom}]$ & Bias \\
\midrule
MolQED       & $0.450$ & $0.476$ & $+0.027$ \\
DiffBases    & $0.197$ & $0.215$ & $+0.018$ \\
TSP-50       & $0.126$ & $0.135$ & $+0.008$ \\
FlatPoly     & $0.077$ & $0.083$ & $+0.006$ \\
TSP-100      & $0.018$ & $0.022$ & $+0.004$ \\
Erd\H{o}s    & $0.476$ & $0.475$ & $-0.001$ \\
MaxCut       & $0.000$ & $0.000$ & $+0.000$ \\
LJ-n41       & $0.000$ & $0.000$ & $+0.000$ \\
CircPack     & $1.000$ & $1.000$ & $+0.000$ \\
\midrule
Mean         &         &         & $+0.006$ \\
\bottomrule
\end{tabular}
\end{table}

Correcting the denominator shifts all protocol AggMEG values by $+0.006$ (uniform, since the correction is task-level). No protocol changes sign. MoA-diverse moves from $-0.022$ to ${\sim}{-}0.016$; its 95\% CI still spans zero in both cases. The bias is conservative; it makes MEG harder to achieve, so all negative findings are robust to correction. The largest per-task bias ($+0.027$ on MolQED) occurs where the three control means are close and noisy, exactly as expected from order-statistic theory.

\paragraph{Anchor invariance of per-task MEG rankings.}
A natural concern is whether the $Q$-normalization anchors ($s_\mathrm{bad}$, $s_\mathrm{ref}$) influence MEG rankings. Within a single task, $Q$ is an affine function of the raw score: $Q(s) = a \cdot s + b$ where the slope $a$ and intercept $b$ depend only on anchors and direction. MEG is a difference of $Q$~values: $\meg = Q(s_p) - \max_c Q(s_c)$. Since $\max_c(a \cdot s_c + b) = a \cdot \min_c(s_c) + b$ for $a < 0$ (minimize tasks; analogously for maximize), the intercept $b$ cancels and $\meg = a \cdot (s_p - s^*_c)$, where $s^*_c$ is the best-control raw score. Perturbing $s_\mathrm{bad}$ or $s_\mathrm{ref}$ changes only $a$, which is a common positive scalar for all protocols, so \textbf{per-task MEG rankings and signs are invariant to anchor choice} (up to the $[0,1]$ clamp, which affects only protocols already at $Q{=}0$ or $Q{=}1$). Across tasks, changing anchors rescales each task's MEG contribution to the aggregate, effectively reweighting tasks. This is the same sensitivity already tested by the leave-one-task-out analysis (Appendix~\ref{app:robustness}), which shows that aggregate rankings are stable within broad tiers but not at the individual-protocol level.

\section{Extended MIG: All Protocol Families}
\label{app:mig}

MIG (Eq.~\ref{eq:mig}) is computed for five multi-agent protocol families in both same-model and diverse-model-family configurations across all nine tasks. For MoA, the no-interact baseline is best-of-proposals from the same three models. For other protocols, it is the best single-agent output from the same model-family set: single-shot for homo protocols; $\max(\text{SS-Claude}, \text{SS-GPT4o}, \text{SS-Gemini})$ per task for mixed protocols. The homo baseline is weaker (one model vs.\ best of three), so homo MIG values are not directly comparable to mixed MIG values, but within each column, the ranking is meaningful.

\begin{table}[h]
\caption{Extended MIG across five protocol families (8 tasks; TSP-50 excluded because MoA-nosynth was not run on that task). Homo: same-model agents, baseline = single-shot. Mixed: diverse model families, baseline = best individual model family. Only MoA-diverse has positive mixed MIG. $\Delta$: mixed $-$ homo MIG.}
\label{tab:mig}
\centering
\small
\begin{tabular}{l r r r}
\toprule
Protocol family & Homo MIG & Mixed MIG & $\Delta$ \\
\midrule
MoA           & $+0.012$ & $\mathbf{+0.016}$ & $+0.004$ \\
MAgICoRe      & $+0.044$ & $-0.035$ & $-0.078$ \\
Debate        & $+0.012$ & $-0.078$ & $-0.089$ \\
HPE           & $-0.163$ & $-0.145$ & $+0.018$ \\
Homo-chain    & $+0.051$ & $-0.023^\dagger$ & $-0.074$ \\
\midrule
Grand mean    & $-0.009$ & $-0.053$ & \\
\bottomrule
\end{tabular}

\smallskip\noindent\small $^\dagger$Cross-chain (mixed variant of homo-chain).
\end{table}

The key pattern: MoA-diverse is the only protocol where mixing diverse model families yields positive MIG ($+0.016$; $\Delta{=}+0.004$). For MAgICoRe and Debate, diversity \emph{reduces} MIG ($\Delta{=}-0.078$, $-0.089$), meaning their interaction steps, adversarial critique and iterative refinement, push diverse agents toward consensus and destroy the complementary information that diversity provides. MoA's one-shot synthesis avoids this because the synthesizer sees all proposals in parallel without iterative convergence pressure.

HPE shows a positive $\Delta$ ($+0.018$), but both homo and mixed MIG are strongly negative, so the improvement reflects less severe failure rather than genuine benefit.

\section{Composability Prediction}
\label{app:composability}

A natural concern is that the diversity finding is post-hoc: we observed that diversity helps on some tasks and rationalized it afterward. As a plausibility check (not strong independent evidence), we define a composability criterion from task structure alone, classify all nine tasks \emph{before} examining the diversity ablation results, and check whether the classification aligns with the observed diversity effects. This check has a significant limitation: the task designer also performed the classification, so design-time intuitions may leak into the labels. The result is suggestive but not confirmatory; independent raters would be needed for stronger validation.

\paragraph{Definition.}
A task is \emph{compositionally separable} if: \textbf{(C1)} a solution can be improved by combining independently generated components (e.g., merging subsets, recombining fragments) without global re-optimization; and \textbf{(C2)} modifying one component does not invalidate the entire solution. Tasks satisfying both conditions are classified HIGH; tasks violating both are LOW.

\begin{table}[h]
\caption{Composability classification (pre-specified, from task structure only) vs.\ observed diversity effect from the $2{\times}2$ ablation. $\Delta_\mathrm{div}$: diverse$-$same-model mean $Q$ (averaged over synthesis conditions). The alignment is 8/9 tasks; MolQED is the only miss.}
\label{tab:composability}
\centering
\small
\begin{tabular}{l l l r l}
\toprule
Task & Composability & Structural reason & $\Delta_\mathrm{div}$ & Prediction \\
\midrule
DiffBases    & HIGH   & Set-union; no coupling       & $+0.13$ & \checkmark \\
Erd\H{o}s    & HIGH   & Regional independence        & $+0.45$ & \checkmark \\
MolQED       & HIGH   & Modular substructure         & $-0.04$ & $\times$ \\
MaxCut       & MEDIUM & Local flips, neighbor-coupled &  $^\dagger$     &   \\
CircPack     & LOW    & Global overlap constraint    &  $^\dagger$     & \checkmark \\
FlatPoly     & LOW    & Global spectral coupling     &  $^\dagger$     & \checkmark \\
TSP-100      & LOW    & Permutation non-decomposable &  $^\dagger$     & \checkmark \\
TSP-50       & LOW    & Permutation non-decomposable &  $-0.01^\ddagger$ & \checkmark \\
LJ-n41       & LOW    & All-pairs energy coupling    &  $^\dagger$     & \checkmark \\
\bottomrule
\end{tabular}

\smallskip\noindent\small $^\dagger$Low-variation tasks: diversity effect not measurable ($Q \approx 0$ for all conditions). $^\ddagger$TSP-50 is discriminative; $\Delta_\mathrm{div}$ computed from MoA-diverse vs.\ MoA-same-model (not from the $2{\times}2$).
\end{table}

\paragraph{Results.}
The composability criterion aligns with the observed pattern for the two tasks where diversity produces large positive effects (DiffBases $+0.13$, Erd\H{o}s $+0.45$), the four LOW tasks where all protocols show low variation ($Q \approx 0$), and TSP-50 (LOW, with a small negative diversity effect consistent with the classification), but misses MolQED (predicted HIGH): despite molecular structure being modular in principle, the diverse models generate chemically incompatible SMILES structures that cannot be usefully composed; the ``modularity'' exists in chemical space, not in the string representation the models actually produce.

\paragraph{Interpretation.}
The composability criterion is a useful but imperfect heuristic. It works when the task's mathematical decomposability aligns with the solution representation the models use (sets for DiffBases, functions for Erd\H{o}s). It fails when domain composability does not map to representational composability (SMILES strings for MolQED). A refined criterion would need to account for the model's solution representation, not just the task's abstract structure. We present this as a partial validation: the alignment is 8/9 by task and 2/3 on the HIGH-composability group where it is most testable, but it is not a complete mechanistic explanation. An out-of-suite probe (Appendix~\ref{app:external}) provides directionally consistent evidence: of three external tasks, only the composable one (Weighted MAX-SAT) shows a positive diversity trend ($+2.4$, CI $[-0.2, +5.4]$); the non-composable tasks (TSP, Number Partition) do not.

\section{Out-of-Suite Probe}
\label{app:external}

To probe generalization beyond \textsc{Veil}, we run the same $2{\times}2$ factorial (diversity $\times$ synthesis, $N{=}3$ proposals, 5 seeds) on three combinatorial optimization tasks outside our benchmark suite: Travelling Salesman (25 cities, non-composable), Weighted MAX-SAT (60 variables, 200 clauses, composable), and Number Partitioning (40 integers, weakly composable). All use local deterministic verifiers with \emph{no} hidden evaluator---dev scores are final scores. MAX-SAT is the only task with high composability: its sparse clause-variable graph creates natural variable clusters that can be optimized independently and recombined.

\begin{table}[h]
\centering
\small
\caption{External $2{\times}2$ probe on three out-of-suite tasks. Values are mean dev scores with bootstrap 95\% CIs ($n{=}5$ seeds, 10{,}000 resamples). Higher is better. The composable task (MAX-SAT) shows a positive diversity trend; the non-composable tasks do not. Partition shows a significant synthesis effect.}
\label{tab:external}
\begin{tabular}{l cc cc cc}
\toprule
 & \multicolumn{2}{c}{\textbf{TSP-25}} & \multicolumn{2}{c}{\textbf{MAX-SAT-60}} & \multicolumn{2}{c}{\textbf{Partition-40}} \\
\cmidrule(lr){2-3}\cmidrule(lr){4-5}\cmidrule(lr){6-7}
 & Same & Diverse & Same & Diverse & Same & Diverse \\
\midrule
No Synth & $\text{-}1218\;{\scriptstyle[\text{-}1256,\text{-}1178]}$ & $\text{-}1181\;{\scriptstyle[\text{-}1234,\text{-}1124]}$ & $1009\;{\scriptstyle[1002,1015]}$ & $1015\;{\scriptstyle[1008,1022]}$ & $\text{-}1585\;{\scriptstyle[\text{-}2074,\text{-}1088]}$ & $\text{-}1292\;{\scriptstyle[\text{-}1813,\text{-}787]}$ \\
Synth    & $\text{-}1145\;{\scriptstyle[\text{-}1203,\text{-}1066]}$ & $\text{-}1174\;{\scriptstyle[\text{-}1210,\text{-}1131]}$ & $1014\;{\scriptstyle[1007,1020]}$ & $1012\;{\scriptstyle[1006,1019]}$ & $\text{-}624\;{\scriptstyle[\text{-}850,\text{-}355]}$  & $\text{-}547\;{\scriptstyle[\text{-}1095,\text{-}244]}$  \\
\midrule
Div.\ $\Delta$ & \multicolumn{2}{c}{$+4\;{\scriptstyle[\text{-}57,+66]}$} & \multicolumn{2}{c}{$+2.4\;{\scriptstyle[\text{-}0.2,+5.4]}$} & \multicolumn{2}{c}{$+189\;{\scriptstyle[\text{-}141,+499]}$} \\
Synth.\ $\Delta$ & \multicolumn{2}{c}{$+40\;{\scriptstyle[\text{-}4,+104]}$} & \multicolumn{2}{c}{$+1.0\;{\scriptstyle[\text{-}1.6,+3.6]}$} & \multicolumn{2}{c}{$\mathbf{+857}\;{\scriptstyle[+341,+1364]}$} \\
\bottomrule
\end{tabular}
\end{table}

\textbf{Composability pattern.} The three tasks span a composability gradient. On TSP (non-composable: sub-tours cannot be merged), the diversity effect is null ($+4$, CI $[-57, +66]$). On MAX-SAT (composable: sparse clause structure allows variable-level recombination), diversity shows a positive trend ($+2.4$, CI $[-0.2, +5.4]$) that is directionally consistent with the composability prediction but \emph{not significant} at $n{=}5$; we report this as suggestive directional evidence, not a confirmed effect. On Partition (weakly composable), the effect is positive but wide ($+189$, CI $[-141, +499]$). The MAX-SAT result also exhibits the diversity tax in miniature: without synthesis, diverse agents outscore same-model agents by $+6.6$; with synthesis, the advantage reverses to $-1.8$, consistent with the aggregation step homogenizing the independent proposals.

\textbf{Synthesis contrast.} On Partition, synthesis shows a large significant effect ($+857$, CI $[+341, +1364]$) under visible-only scoring. On the main benchmark under hidden evaluation, synthesis is non-significant ($-0.01$, CI $[-0.11, +0.09]$). These are different tasks under different evaluation regimes, so the comparison is suggestive rather than causal: the pattern is \emph{consistent with} hidden evaluation suppressing scorer-specific synthesis gains, but could also reflect task-level differences in synthesis value.

\section{Budget Scaling}
\label{app:budget-scaling}

To test whether the negative MEG finding is an artifact of insufficient budget, we re-run the four single-agent controls and MoA-diverse at $2{\times}$ budget ($T{=}400$K tokens, $W{=}1200$s, $C{=}60$s CPU, $K{=}50$ eval calls) on the three discriminative tasks (DiffBases, Erd\H{o}s, MolQED; 5 seeds each; $N{=}74$ cells, one BoN cell timed out).

\begin{table}[h]
\caption{Budget scaling: mean $Q$ at $1{\times}$ vs.\ $2{\times}$ budget on three discriminative tasks. Controls (BoN, VGS) improve more than MoA-diverse on aggregate. MoA-diverse trails the max control on all three tasks at $2{\times}$.}
\label{tab:budget2x}
\centering
\small
\begin{tabular}{l r r r r r r}
\toprule
& \multicolumn{2}{c}{DiffBases} & \multicolumn{2}{c}{Erd\H{o}s} & \multicolumn{2}{c}{MolQED} \\
\cmidrule(lr){2-3}\cmidrule(lr){4-5}\cmidrule(lr){6-7}
Protocol & $1{\times}$ & $2{\times}$ & $1{\times}$ & $2{\times}$ & $1{\times}$ & $2{\times}$ \\
\midrule
Single-shot & $0.00$ & $0.00$ & $0.00$ & $0.00$ & $0.34$ & $0.17$ \\
Self-Refine & $0.09$ & $0.16$ & $0.48$ & $0.31$ & $0.40$ & $0.46$ \\
Best-of-$N$ & $0.20$ & $0.13$ & $0.11$ & $\mathbf{0.68}$ & $0.45$ & $0.52$ \\
VGS         & $0.27$ & $\mathbf{0.39}$ & $0.00$ & $0.38$ & $0.41$ & $0.47$ \\
\midrule
Max control & $0.27$ & $0.39$ & $0.48$ & $0.68$ & $0.45$ & $0.52$ \\
\midrule
MoA-diverse & $0.29$ & $0.34$ & $0.50$ & $0.57$ & $0.44$ & $0.39$ \\
\bottomrule
\end{tabular}
\end{table}

The max-control denominator increases on all three tasks ($+0.13$, $+0.21$, $+0.07$), meaning the MEG bar rises with budget. VGS benefits from more evaluator calls (Erd\H{o}s: $0.00 \to 0.38$; DiffBases: $0.27 \to 0.39$), while Best-of-$N$ gains from more independent samples (Erd\H{o}s: $0.11 \to 0.68$).

MoA-diverse also improves on DiffBases ($0.29 \to 0.34$) and Erd\H{o}s ($0.50 \to 0.57$), but drops on MolQED ($0.44 \to 0.39$). Critically, MoA-diverse trails the $2{\times}$ max control on all three tasks: by $-0.05$ on DiffBases, $-0.11$ on Erd\H{o}s, and $-0.13$ on MolQED. The controls scale faster than MoA-diverse, widening the MEG gap rather than closing it. The negative MEG finding is therefore not a budget artifact: even MoA-diverse, the protocol most favored by the $1{\times}$ analysis, does not catch the single-agent controls at $2{\times}$ budget.

\section{Robustness Checks}
\label{app:robustness}

\paragraph{$2{\times}2$ linear model.}
We fit $Q \sim \text{diversity} + \text{synthesis} + \text{diversity}{\times}\text{synthesis} + \text{task dummies}$ via OLS on the balanced $N{=}120$ observations from the three high-variation tasks (DiffBases, Erd\H{o}s, MolQED) with 4 matched protocol arms $\times$ 10 seeds each. Results:

\begin{center}
\small
\begin{tabular}{l r r}
\toprule
Coefficient & Estimate & 95\% CI \\
\midrule
Diversity       & $+0.188$ & $[+0.064, +0.309]$ \\
Synthesis       & $-0.010$ & $[-0.110, +0.091]$ \\
Diversity$\times$Synthesis & $+0.046$ & $[-0.131, +0.221]$ \\
\midrule
Div $-$ Synth   & $+0.198$ & $[+0.071, +0.324]$ \\
\bottomrule
\end{tabular}
\end{center}

\noindent $R^{2}{=}0.21$, $N{=}120$. CIs via 10{,}000 percentile bootstrap. Diversity individually excludes zero; synthesis and the interaction do not. Diversity exceeds synthesis in 99.9\% of resamples.

\paragraph{Random-intercept check.}
A task-random-intercept REML fit on the same data gives the same fixed-effect estimates: diversity $+0.188$ \,[+0.064, +0.313], synthesis $-0.010$ \,[-0.134, +0.115], interaction $+0.046$ \,[-0.130, +0.222]. Task-level intercept variance is small ($\hat{\tau}^2 \approx 0.0049$), which is consistent with the fixed-effects model and indicates that the diversity signal is not driven by a single task intercept.

\paragraph{Leave-one-task-out.}
Dropping each of the 9 tasks and recomputing aggregate MEG over the remaining 8 reveals sensitivity to task composition. Removing CircPack shifts MoA to rank~1 with positive AggMEG ($+0.005$); removing DiffBases drops MoA to $-0.036$ (rank~2). Rankings are stable only within broad tiers (negative-near-zero vs.\ strongly-negative), not at the individual-protocol level.

\paragraph{Median aggregation.}
Replacing mean with median across tasks changes the top ranking: Best-of-$N$ and VGS tie at median MEG${=}0.000$ (ranks~1--2), MoA drops to rank~3, and Debate jumps from rank~8 to rank~4. The mean ranking is driven partly by a few high-variation tasks (Erd\H{o}s, CircPack).

\paragraph{Jackknife standard errors.}
Leave-one-seed-out jackknife SE for aggregate MEG:

\begin{center}
\small
\begin{tabular}{l r r r}
\toprule
Protocol & AggMEG & Jack SE & $\pm2$SE interval \\
\midrule
MoA         & $-0.022$ & $0.030$ & $[-0.082, +0.038]$ \\
Self-Refine & $-0.038$ & $0.013$ & $[-0.064, -0.012]$ \\
Best-of-$N$ & $-0.051$ & $0.041$ & $[-0.133, +0.031]$ \\
\bottomrule
\end{tabular}
\end{center}

\noindent MoA and Best-of-$N$ are the only protocols whose $\pm2$SE intervals span zero; Best-of-$N$'s wide interval reflects sensitivity to the Erd\H{o}s task (where it is the binding control, producing MEG $= 0$ on some drop-one subsets). Self-Refine excludes zero (negative side). Fine-grained ranking among the top three is not statistically resolved at the current seed count.

\subsection{Multi-Synthesizer Robustness Check}
\label{app:multisynth}

The primary $2{\times}2$ ablation uses Claude Sonnet~4 as the synthesizer in both synthesis cells. To rule out the possibility that the diversity advantage is an artifact of Claude's synthesis, we re-run the synthesis cells with GPT-4o and Gemini~2.5 Flash as the aggregator model, keeping the proposer backbones identical. Table~\ref{tab:multisynth} reports hidden-score $Q$ on DiffBases (the task with the largest diversity effect in the original ablation).

\begin{table}[h]
\caption{Multi-synthesizer check on DiffBases ($Q$-normalized hidden scores, 10 seeds per cell). The diversity advantage is positive for all three synthesizer models.}
\label{tab:multisynth}
\centering
\small
\begin{tabular}{l c c c}
\toprule
Synthesizer & Diverse $\bar{Q}$ & Same $\bar{Q}$ & $\Delta$ \\
\midrule
Claude Sonnet~4 & $0.287$ & $0.070$ & $+0.217$ \\
GPT-4o          & $0.257$ & $0.118$ & $+0.139$ \\
Gemini 2.5 Flash & $0.402$ & $0.230$ & $+0.172$ \\
\bottomrule
\end{tabular}
\end{table}

\noindent The diversity advantage is consistent in sign across all three synthesizer models ($\Delta > 0$ in every row), ranging from $+0.14$ to $+0.22$. The synthesizer model affects overall performance level (Gemini produces the highest $\bar{Q}$ in both conditions) but does not eliminate the diversity effect. This confirms that the primary finding---model-family diversity among proposers matters more than the aggregation mechanism---is not confounded with the synthesizer identity.

\paragraph{Missing data.}
Twelve of 535 expected core-protocol result files (2.2\%) are missing, concentrated in TSP-100 controls (Self-Refine: 1/5 seeds, Best-of-$N$: 3/5, VGS: 4/5) and HPE on DiffBases and Erd\H{o}s. TSP-100 is a low-spread task ($Q \approx 0$ for all protocols), so the missing control seeds have minimal effect on aggregate MEG. HPE's sparse data is noted in the main tables.

\section{Extending Veil}
\label{app:extensibility}

\textsc{Veil} is designed to grow. A new task enters the benchmark if it satisfies four requirements:

\begin{enumerate}
\item \textbf{Deterministic verifier.} The task must have a programmatic scoring function that maps a solution artifact to a real-valued score without stochasticity. Human or LLM-based evaluation is out of scope.
\item \textbf{Hidden-evaluator split.} Two verifier variants must exist: a \emph{development} evaluator that agents observe during search, and a \emph{hidden} evaluator that differs in input, scoring components, or both. The hidden variant must be constructible from the task's domain structure---e.g., perturbing problem instances (Type~II) or increasing coverage requirements (Type~I).
\item \textbf{Non-trivial search headroom.} Under a baseline single-shot call, the task must produce non-saturated $Q$ scores (i.e., $Q < 0.9$) so that iterative protocols have room to improve. Tasks where all protocols score $Q \approx 0$ are informative as capability-frontier markers but do not contribute to protocol discrimination.
\item \textbf{Frozen normalization anchors.} Two reference scores ($s_\mathrm{bad}$, $s_\mathrm{ref}$) must be fixed before any protocol comparison begins, following the procedure in Appendix~\ref{app:normalization}.
\end{enumerate}

\paragraph{Hidden-evaluator design patterns.}
Two patterns cover all current tasks. \textbf{Type~I (input replacement)}: the hidden evaluator uses the same scoring function but on a harder input instance---e.g., the full graph instead of the 80\%-edge training subgraph (MaxCut), or a stricter tolerance (Circle Packing). This tests whether the solution generalizes beyond the visible instance. \textbf{Type~II (criterion addition)}: the hidden evaluator adds a scoring component invisible during search---e.g., perturbed city coordinates (TSP), adversarial evaluation points (Flat Polynomials), or a synthesizability penalty (MolQED). This tests whether the agent optimizes the true objective or a visible proxy.

\paragraph{Adding a task in practice.}
Implement \texttt{devVerifier(artifact) $\to$ float} and \texttt{hiddenVerifier(artifact) $\to$ float}, define the solution schema, set normalization anchors from pilot runs, and register the task in the configuration. The budget vector ($T, W, C, K$) is shared across all tasks; per-task adjustment is not needed. Code templates and a worked example (adding a new combinatorial task) will be released with the benchmark.

\section{Protocol Reimplementation Fidelity}
\label{app:reimplementation}

All multi-agent protocols in \textsc{Veil} are our reimplementations based on published descriptions. We have not validated against the original authors' codebases. This section documents the relationship between our implementations and the source papers.

\begin{table}[h]
\caption{Protocol reimplementation summary. ``Original domain'' is the evaluation setting used in the source paper. None of the original papers evaluate on structured optimization with hidden verifiers, so direct numerical comparison is not possible.}
\label{tab:reimplementation}
\centering
\small
\begin{tabular}{l l l l}
\toprule
Protocol & Source & Original domain & Key adaptation \\
\midrule
Self-Refine & \citet{madaan2023selfrefine} & Code, reasoning & Verifier feedback replaces self-critique \\
Debate & \citet{du2024debate} & Factuality, reasoning & Artifact generation replaces answer selection \\
MAgICoRe & \citet{chen2025magicore} & Math reasoning & Solver/Reviewer/Refiner over artifacts \\
MoA & \citet{wang2025moa} & AlpacaEval, MT-Bench & Proposers generate artifacts, not text \\
HPE & \citet{wang2023plansolve} & Math word problems & Planner decomposes; executors produce artifacts \\
\bottomrule
\end{tabular}
\end{table}

\paragraph{Domain mismatch.}
The original papers evaluate on open-ended language tasks (factuality, reasoning, chat quality) using LLM judges or exact-match correctness. \textsc{Veil} evaluates on structured optimization tasks with deterministic hidden verifiers. Because the task domains do not overlap, we cannot directly compare our scores against published baselines. The negative MEG finding should therefore be interpreted as specific to structured optimization under hidden evaluation, not as a rebuttal of the original papers' claims in their respective domains.

\paragraph{Structural fidelity.}
For each protocol, we preserve the core interaction pattern: Self-Refine's iterative loop with scorer feedback, debate's adversarial exchange between two agents plus synthesis, MAgICoRe's coarse-to-fine Solver$\to$Reviewer$\to$Refiner chain, MoA's parallel proposers plus one-shot synthesis, and HPE's planner-executor decomposition. The primary adaptation is output format: all protocols generate structured solution artifacts (e.g., integer sets, coordinate arrays, SMILES strings) rather than natural-language answers. Our Self-Refine variant replaces the self-critique step of \citet{madaan2023selfrefine} with direct verifier feedback, making it a stronger control baseline; this deviation is noted in Limitation~(viii).

\begin{thebibliography}{99}

\bibitem[Bickerton et~al.(2012)]{bickerton2012qed}
Bickerton, G.~R., et~al. (2012).
Quantifying the chemical beauty of drugs.
\emph{Nature Chemistry}, 4(2), 90--98.

\bibitem[Ertl and Schuffenhauer(2009)]{ertl2009sascore}
Ertl, P. and Schuffenhauer, A. (2009).
Estimation of synthetic accessibility score of drug-like molecules based on molecular complexity and fragment contributions.
\emph{Journal of Cheminformatics}, 1(1), 8.

\bibitem[Chen et~al.(2025)]{chen2025magicore}
Chen, J.~C.-Y., et~al. (2025).
MAgICoRe: Multi-agent, iterative, coarse-to-fine refinement for reasoning.
EMNLP, arXiv:2409.12147.

\bibitem[Du et~al.(2024)]{du2024debate}
Du, Y., et~al. (2024).
Improving factuality and reasoning in language models through multiagent debate.
ICML, arXiv:2305.14325.

\bibitem[Madaan et~al.(2023)]{madaan2023selfrefine}
Madaan, A., et~al. (2023).
Self-Refine: Iterative refinement with self-feedback.
NeurIPS, arXiv:2303.17651.

\bibitem[Romera-Paredes et~al.(2024)]{romeraparedes2024funsearch}
Romera-Paredes, B., et~al. (2024).
Mathematical discoveries from program search with large language models.
Nature, 625, 468--475.

\bibitem[Yang et~al.(2024)]{yang2024opro}
Yang, C., et~al. (2024).
Large language models as optimizers.
ICLR 2024, arXiv:2309.03409.

\bibitem[Wang et~al.(2025)]{wang2025moa}
Wang, J., et~al. (2025).
Mixture-of-Agents enhances large language model capabilities.
ICLR Spotlight, arXiv:2406.04692.

\bibitem[Snell et~al.(2024)]{snell2024scaling}
Snell, C., Lee, J., Xu, K., and Kumar, A. (2024).
Scaling LLM test-time compute optimally can be more effective than scaling model parameters.
arXiv:2408.03314.

\bibitem[Xu et~al.(2026)]{xu2026single}
Xu, J., et~al. (2026).
Rethinking the value of multi-agent workflow: A strong single agent baseline.
arXiv:2601.12307.

\bibitem[Tran and Kiela(2026)]{tran2026equal}
Tran, D. and Kiela, D. (2026).
Single-Agent LLMs Outperform Multi-Agent Systems on Multi-Hop Reasoning Under Equal Thinking Token Budgets.
arXiv:2604.02460.

\bibitem[Novikov et~al.(2025)]{novikov2025alphaevolve}
Novikov, A., et~al. (2025).
AlphaEvolve: A coding agent for scientific and algorithmic discovery.
arXiv:2506.13131.

\bibitem[Georgiev et~al.(2025)]{georgiev2025alphaevolve}
Georgiev, B., G\'omez-Serrano, J., Tao, T., and Wagner, A.~Z. (2025).
Mathematical exploration and discovery at scale.
arXiv:2511.02864.

\bibitem[Li et~al.(2023)]{li2023camel}
Li, G., et~al. (2023).
CAMEL: Communicative agents for ``mind'' exploration of large language model society.
NeurIPS, arXiv:2303.17760.

\bibitem[Wu et~al.(2023)]{wu2023autogen}
Wu, Q., et~al. (2023).
AutoGen: Enabling next-gen LLM applications via multi-agent conversation.
arXiv:2308.08155.

\bibitem[Chen et~al.(2023)]{chen2023agentverse}
Chen, W., et~al. (2023).
AgentVerse: Facilitating multi-agent collaboration and exploring emergent behaviors.
arXiv:2308.10848.

\bibitem[Hong et~al.(2023)]{hong2023metagpt}
Hong, S., et~al. (2023).
MetaGPT: Meta programming for a multi-agent collaborative framework.
arXiv:2308.00352.

\bibitem[Wang et~al.(2023)]{wang2023plansolve}
Wang, L., et~al. (2023).
Plan-and-solve prompting: Improving zero-shot chain-of-thought reasoning by large language models.
ACL 2023, arXiv:2305.04091.

\bibitem[Krogh and Vedelsby(1994)]{krogh1994neural}
Krogh, A. and Vedelsby, J. (1994).
Neural network ensembles, cross validation, and active learning.
\emph{Advances in Neural Information Processing Systems}, 7.

\bibitem[Lorenz et~al.(2011)]{lorenz2011social}
Lorenz, J., Rauhut, H., Schweitzer, F., and Helbing, D. (2011).
How social influence can undermine the wisdom of crowd effect.
\emph{Proceedings of the National Academy of Sciences}, 108(22):9020--9025.

\end{thebibliography}

\end{document}
